% Manuscript fragment. Requires amsmath, amssymb, amsthm.
\subsection{Explicit guarded normalization of record readability}
\label{sec:guarded-provenance}

The path and readability definitions above already account for copying and
later key disclosure. We now give their explicit guarded normalization and
an executable realization for a record-decryption specialization. Here key
derivation has fixed public edges independent of recovered plaintext. The
optional plaintext-dependent key-extraction extension remains as defined
above and is outside this specialization. The set $\Omega_{\rm src}$ below
contains reservation-source identities; it is distinct from saved record
versions and from the original-payload table $\mathcal O[v]$. Likewise,
$\mathcal E^{\rm src}_t$ projects derivability onto protected source values,
whereas $\mathcal C_t$ and $\mathcal L_t$ track acquired and readable versions.
The graph is a fixed representation of admitted retained objects, not a
replacement for the provenance transition rules or the byte renderer.


A source identifier names one original sampling operation.  An acquisition
history records what the adversary obtains, including objects obtained at
different times.  These are distinct objects of the model: a copy preserves
the former, while an additional acquisition can change the latter.

\paragraph{Immutable record graph.}
Let $\mathcal R$ be a finite set of immutable encrypted record objects.
Each $r\in\mathcal R$ has a matching decryption-key identity $\kappa(r)$,
a set of embedded child records $\operatorname{child}(r)$, and a set
$\operatorname{orig}(r)$ of protected source values in its plaintext.
Containment is acyclic.  Record identities distinguish ciphertext objects;
source identities distinguish original sampling operations.  Several
records can contain the same source.  A rewrapped record has a new record
identity and may have a new key, while retaining its source identity.
Key identities are equivalence classes of equal actual keys. Equal key bytes
therefore have one identity across separate API objects.

Let $\mathcal D\subseteq\mathcal K\times\mathcal K$ be fixed public key
derivation edges.  An edge $(k,k')$ means that possession of $k$ permits
deterministic computation of $k'$.  This interface represents, for example,
derivation of an operation-wrapping key from the journal key.  In the
theorem below, these edges do not depend on plaintext recovered from records.
Write $D^*$ for their reflexive transitive closure.

Let $\Omega_{\rm src}=\bigcup_{r\in\mathcal R}\operatorname{orig}(r)$ be the
declared source domain.  At a trace prefix $t$, let $A_t\subseteq\mathcal R$ be all record objects
ever acquired by that prefix, $K_t\subseteq\mathcal K$ all directly
disclosed keys, and $\mathcal S^{\rm dir}_t\subseteq\Omega_{\rm src}$ all directly disclosed source values, represented
by their source identities.  These are cumulative sets.  The decryption-key
closure is
\[
  K_t^*=\{k':\exists k\in K_t,\ (k,k')\in D^*\}.
\]
The concrete derivation rules are
\[
 \frac{k\quad(k,k')\in\mathcal D}{k'},\qquad
 \frac{r\quad\kappa(r)}{c}\ (c\in\operatorname{child}(r)),\qquad
 \frac{r\quad\kappa(r)}{o}\ (o\in\operatorname{orig}(r)).
\]
Their least closure starts from $A_t,K_t,\mathcal S^{\rm dir}_t$.  A fact $o$ in the source
sort means that the protected value of source $o$ is derivable; it does
not mean that a public identifier has become known.  The source projection
of this closure is denoted $\mathcal E^{\rm src}_t$.

\paragraph{Guarded routes.}
For every path $p=(r_0,\ldots,r_\ell)$ in the containment graph ending
at a record with $o\in\operatorname{orig}(r_\ell)$, define
\[
 g(p)=\{\kappa(r_0),\ldots,\kappa(r_\ell)\}.
\]
The normalized table contains a route $(r_0,g(p),o,p)$ for every such
path.  Its exposure expression is
\[
 \widehat{\mathcal E}^{\rm src}_t=\mathcal S^{\rm dir}_t\ \cup\
 \{o:\exists(r,g,o,p),\ r\in A_t\ \land\ g\subseteq K_t^*\}.
 \tag{GP}
\]
Alternative records and paths are disjoined, while all keys on one path
are conjoined.  Equal key identities occur only once in the guard set.

\begin{theorem}[Exact guarded normalization]
\label{thm:guarded-normalization}
For every finite record graph and derivation interface above, and every
acquisition prefix $(A_t,K_t,\mathcal S^{\rm dir}_t)$, the concrete and normalized exposure
sets agree: $\mathcal E^{\rm src}_t=\widehat{\mathcal E}^{\rm src}_t$.  The equality applies to all prefixes of
an adaptively chosen history.
\end{theorem}
\begin{proof}
Key derivations depend only on disclosed keys and the fixed relation
$\mathcal D$, so their least closure is exactly $K_t^*$.
For the forward inclusion, take a source $o\in \mathcal E^{\rm src}_t$.
If it was directly disclosed then $o\in \mathcal S^{\rm dir}_t$.
Otherwise its first derivation decrypts a record $r_\ell$ with the
matching key.  If that record was not acquired directly, its first
availability follows by decrypting a parent.  Follow these first
derivations backwards.  The finite acyclic record graph makes this
process terminate at an acquired record $r_0$.  Every intervening key
lies in $K_t^*$, yielding a route satisfying (GP).
For the reverse inclusion, take a satisfied route.  Starting with its
acquired first record, successively decrypt along the path.  Every
required key is in $K_t^*$, so each child and finally $o$ is in the
concrete least closure.  Direct disclosures are in that closure by
definition.  Both inclusions hold for each fixed prefix; no independence
or advance choice of the acquisition schedule was used.
\end{proof}

\paragraph{Efficient evaluation.}
Equation (GP) is a declarative normal form: a compressed record DAG can
have exponentially many paths.  An efficient leakage evaluator instead
computes the iterative decryption and key-derivation fixed point, without
enumerating those paths.  Let $N$ be the total number of record, key and
source atoms, and let $M$ count rule and payload incidences.  There are at
most $N$ strict-growth rounds; repeated rule scans therefore take
$O(N(N+M))$ time and $O(N+M)$ space under standard indexed set and
dictionary representations, with source-dictionary operations included
in this accounting.  The finite checker materializes the route table only
for its declared small fixtures, to compare that normal form against the
forward evaluator.

\begin{corollary}[Copies, deletion and late disclosure]
\label{cor:guarded-history}
Exposure is monotone under additional acquisitions, direct disclosures
and key disclosures.  Copying or rewrapping a source adds routes to the
same source identity.  Local deletion of a record previously acquired
does not remove its routes from adversarial knowledge.  In particular,
a later matching-key disclosure can expose a source in a ciphertext
acquired before that disclosure and deleted locally in the meantime.
\end{corollary}
\begin{proof}
The cumulative sets and their key closure only grow, and every predicate
in (GP) is positive.  A source-preserving copy changes a route's record
identity or key guards but not its terminal source.  Local deletion
restricts future acquisition from that local store; it removes no member
of $A_t$ and therefore cannot invalidate a satisfied route.
\end{proof}

\paragraph{Interface significance.}
The theorem gives an exact finite normalization of the chosen
record-decryption semantics.  It supports a precise original-operation
leakage oracle: a query for source $o$ can be evaluated against the
complete acquisition history, with a path and its key guards as a witness.
The cryptographic transition from native ciphertexts and public
commitments to this interface has a separate simulation obligation.
The derivation rules do not model guessing, cryptanalysis, timing leakage
or a quantum decryption oracle.  Accordingly, absence from $\mathcal E^{\rm src}_t$ is
recorded as non-derivability by these rules; message indistinguishability
outside an exposure set is supplied by a security theorem for the
surrounding protocol.

\paragraph{Native shape and finite checks.}
At the fixed research source, operation reservations are encrypted under
a wrapping key derived from the journal key, and tokens are retained
within encrypted durable images or pending targets.  The corresponding
route therefore includes the image key and the token key, with their
actual derivation edge.  A fresh aggregate image can preserve an older
token and its original source.  Public image-head digests are distinct
observations from the record objects needed by the decryption rules.
The executable supplement independently evaluates iterative decryption
and guarded-route normalization, exhausts all acquisition/key/direct-source
subsets for its declared finite graphs, and checks histories with late
key disclosure and local deletion.  Its counts and source hash are in
the accompanying machine-readable result; those checks validate the
implemented finite model and supplement the general proof above.
